% Options for packages loaded elsewhere
\PassOptionsToPackage{unicode}{hyperref}
\PassOptionsToPackage{hyphens}{url}
\documentclass[
]{book}
\usepackage{xcolor}
\usepackage{amsmath,amssymb}
\setcounter{secnumdepth}{5}
\usepackage{iftex}
\ifPDFTeX
  \usepackage[T1]{fontenc}
  \usepackage[utf8]{inputenc}
  \usepackage{textcomp} % provide euro and other symbols
\else % if luatex or xetex
  \usepackage{unicode-math} % this also loads fontspec
  \defaultfontfeatures{Scale=MatchLowercase}
  \defaultfontfeatures[\rmfamily]{Ligatures=TeX,Scale=1}
\fi
\usepackage{lmodern}
\ifPDFTeX\else
  % xetex/luatex font selection
\fi
% Use upquote if available, for straight quotes in verbatim environments
\IfFileExists{upquote.sty}{\usepackage{upquote}}{}
\IfFileExists{microtype.sty}{% use microtype if available
  \usepackage[]{microtype}
  \UseMicrotypeSet[protrusion]{basicmath} % disable protrusion for tt fonts
}{}
\makeatletter
\@ifundefined{KOMAClassName}{% if non-KOMA class
  \IfFileExists{parskip.sty}{%
    \usepackage{parskip}
  }{% else
    \setlength{\parindent}{0pt}
    \setlength{\parskip}{6pt plus 2pt minus 1pt}}
}{% if KOMA class
  \KOMAoptions{parskip=half}}
\makeatother
\usepackage{color}
\usepackage{fancyvrb}
\newcommand{\VerbBar}{|}
\newcommand{\VERB}{\Verb[commandchars=\\\{\}]}
\DefineVerbatimEnvironment{Highlighting}{Verbatim}{commandchars=\\\{\}}
% Add ',fontsize=\small' for more characters per line
\usepackage{framed}
\definecolor{shadecolor}{RGB}{248,248,248}
\newenvironment{Shaded}{\begin{snugshade}}{\end{snugshade}}
\newcommand{\AlertTok}[1]{\textcolor[rgb]{0.94,0.16,0.16}{#1}}
\newcommand{\AnnotationTok}[1]{\textcolor[rgb]{0.56,0.35,0.01}{\textbf{\textit{#1}}}}
\newcommand{\AttributeTok}[1]{\textcolor[rgb]{0.13,0.29,0.53}{#1}}
\newcommand{\BaseNTok}[1]{\textcolor[rgb]{0.00,0.00,0.81}{#1}}
\newcommand{\BuiltInTok}[1]{#1}
\newcommand{\CharTok}[1]{\textcolor[rgb]{0.31,0.60,0.02}{#1}}
\newcommand{\CommentTok}[1]{\textcolor[rgb]{0.56,0.35,0.01}{\textit{#1}}}
\newcommand{\CommentVarTok}[1]{\textcolor[rgb]{0.56,0.35,0.01}{\textbf{\textit{#1}}}}
\newcommand{\ConstantTok}[1]{\textcolor[rgb]{0.56,0.35,0.01}{#1}}
\newcommand{\ControlFlowTok}[1]{\textcolor[rgb]{0.13,0.29,0.53}{\textbf{#1}}}
\newcommand{\DataTypeTok}[1]{\textcolor[rgb]{0.13,0.29,0.53}{#1}}
\newcommand{\DecValTok}[1]{\textcolor[rgb]{0.00,0.00,0.81}{#1}}
\newcommand{\DocumentationTok}[1]{\textcolor[rgb]{0.56,0.35,0.01}{\textbf{\textit{#1}}}}
\newcommand{\ErrorTok}[1]{\textcolor[rgb]{0.64,0.00,0.00}{\textbf{#1}}}
\newcommand{\ExtensionTok}[1]{#1}
\newcommand{\FloatTok}[1]{\textcolor[rgb]{0.00,0.00,0.81}{#1}}
\newcommand{\FunctionTok}[1]{\textcolor[rgb]{0.13,0.29,0.53}{\textbf{#1}}}
\newcommand{\ImportTok}[1]{#1}
\newcommand{\InformationTok}[1]{\textcolor[rgb]{0.56,0.35,0.01}{\textbf{\textit{#1}}}}
\newcommand{\KeywordTok}[1]{\textcolor[rgb]{0.13,0.29,0.53}{\textbf{#1}}}
\newcommand{\NormalTok}[1]{#1}
\newcommand{\OperatorTok}[1]{\textcolor[rgb]{0.81,0.36,0.00}{\textbf{#1}}}
\newcommand{\OtherTok}[1]{\textcolor[rgb]{0.56,0.35,0.01}{#1}}
\newcommand{\PreprocessorTok}[1]{\textcolor[rgb]{0.56,0.35,0.01}{\textit{#1}}}
\newcommand{\RegionMarkerTok}[1]{#1}
\newcommand{\SpecialCharTok}[1]{\textcolor[rgb]{0.81,0.36,0.00}{\textbf{#1}}}
\newcommand{\SpecialStringTok}[1]{\textcolor[rgb]{0.31,0.60,0.02}{#1}}
\newcommand{\StringTok}[1]{\textcolor[rgb]{0.31,0.60,0.02}{#1}}
\newcommand{\VariableTok}[1]{\textcolor[rgb]{0.00,0.00,0.00}{#1}}
\newcommand{\VerbatimStringTok}[1]{\textcolor[rgb]{0.31,0.60,0.02}{#1}}
\newcommand{\WarningTok}[1]{\textcolor[rgb]{0.56,0.35,0.01}{\textbf{\textit{#1}}}}
\usepackage{longtable,booktabs,array}
\usepackage{calc} % for calculating minipage widths
% Correct order of tables after \paragraph or \subparagraph
\usepackage{etoolbox}
\makeatletter
\patchcmd\longtable{\par}{\if@noskipsec\mbox{}\fi\par}{}{}
\makeatother
% Allow footnotes in longtable head/foot
\IfFileExists{footnotehyper.sty}{\usepackage{footnotehyper}}{\usepackage{footnote}}
\makesavenoteenv{longtable}
\usepackage{graphicx}
\makeatletter
\newsavebox\pandoc@box
\newcommand*\pandocbounded[1]{% scales image to fit in text height/width
  \sbox\pandoc@box{#1}%
  \Gscale@div\@tempa{\textheight}{\dimexpr\ht\pandoc@box+\dp\pandoc@box\relax}%
  \Gscale@div\@tempb{\linewidth}{\wd\pandoc@box}%
  \ifdim\@tempb\p@<\@tempa\p@\let\@tempa\@tempb\fi% select the smaller of both
  \ifdim\@tempa\p@<\p@\scalebox{\@tempa}{\usebox\pandoc@box}%
  \else\usebox{\pandoc@box}%
  \fi%
}
% Set default figure placement to htbp
\def\fps@figure{htbp}
\makeatother
\setlength{\emergencystretch}{3em} % prevent overfull lines
\providecommand{\tightlist}{%
  \setlength{\itemsep}{0pt}\setlength{\parskip}{0pt}}
\usepackage[]{natbib}
\bibliographystyle{plainnat}
\usepackage{booktabs}
\usepackage{bookmark}
\IfFileExists{xurl.sty}{\usepackage{xurl}}{} % add URL line breaks if available
\urlstyle{same}
\hypersetup{
  pdftitle={GTI - Gestão e Análise de Dados},
  pdfauthor={Professor Miguél Suares},
  hidelinks,
  pdfcreator={LaTeX via pandoc}}

\title{GTI - Gestão e Análise de Dados}
\author{Professor Miguél Suares}
\date{2026-09-15}

\begin{document}
\maketitle

{
\setcounter{tocdepth}{1}
\tableofcontents
}
\chapter*{Sobre estas anotações}\label{sobre-estas-anotauxe7uxf5es}
\addcontentsline{toc}{chapter}{Sobre estas anotações}

---------------------------------------------------------------------------------------------------------------------------------------

\section*{Calendário das aulas}\label{calenduxe1rio-das-aulas}
\addcontentsline{toc}{section}{Calendário das aulas}

---------------------------------------------------------------------------------------------------------------------------------------

\paragraph{AGOSTO DE 2026}\label{agosto-de-2026}

\begin{longtable}[]{@{}
  >{\raggedright\arraybackslash}p{(\linewidth - 6\tabcolsep) * \real{0.0694}}
  >{\raggedright\arraybackslash}p{(\linewidth - 6\tabcolsep) * \real{0.0903}}
  >{\raggedright\arraybackslash}p{(\linewidth - 6\tabcolsep) * \real{0.1806}}
  >{\raggedright\arraybackslash}p{(\linewidth - 6\tabcolsep) * \real{0.6597}}@{}}
\toprule\noalign{}
\begin{minipage}[b]{\linewidth}\raggedright
Data
\end{minipage} & \begin{minipage}[b]{\linewidth}\raggedright
Dia da Semana
\end{minipage} & \begin{minipage}[b]{\linewidth}\raggedright
Aulas
\end{minipage} & \begin{minipage}[b]{\linewidth}\raggedright
Conteúdo
\end{minipage} \\
\midrule\noalign{}
\endhead
\bottomrule\noalign{}
\endlastfoot
04/08/2026 & Terça-Feira & Apresentação da disciplina & Apresentação do plano de ensino, objetivos, metodologia, avaliações e visão geral da disciplina \\
11/08/2026 & Terça-Feira & Aula 01 & Dados como ativo estratégico e Introdução ao Machine Learning \\
18/08/2026 & Terça-Feira & Aula 02 & Governança de Dados e ciclo de vida de Machine Learning \\
25/08/2026 & Terça-Feira & Aula 03 & Coleta de dados e Regressão Linear \\
\end{longtable}

\paragraph{SETEMBRO DE 2026}\label{setembro-de-2026}

\begin{longtable}[]{@{}
  >{\raggedright\arraybackslash}p{(\linewidth - 6\tabcolsep) * \real{0.1075}}
  >{\raggedright\arraybackslash}p{(\linewidth - 6\tabcolsep) * \real{0.1398}}
  >{\raggedright\arraybackslash}p{(\linewidth - 6\tabcolsep) * \real{0.0753}}
  >{\raggedright\arraybackslash}p{(\linewidth - 6\tabcolsep) * \real{0.6774}}@{}}
\toprule\noalign{}
\begin{minipage}[b]{\linewidth}\raggedright
Data
\end{minipage} & \begin{minipage}[b]{\linewidth}\raggedright
Dia da Semana
\end{minipage} & \begin{minipage}[b]{\linewidth}\raggedright
Aulas
\end{minipage} & \begin{minipage}[b]{\linewidth}\raggedright
Conteúdo
\end{minipage} \\
\midrule\noalign{}
\endhead
\bottomrule\noalign{}
\endlastfoot
01/09/2026 & Terça-Feira & Aula 04 & Organização e preparação dos dados, Regressão Logística e KNN \\
08/09/2026 & Terça-Feira & Aula 05 & Análise de dados, dilema Viés-Variância e Métricas de Avaliação \\
15/09/2026 & Terça-Feira & Aula 06 & Relatórios, KPIs, avaliação de modelos e Árvores de Decisão \\
22/09/2026 & Terça-Feira & NP1 & PROVA \\
29/09/2026 & Terça-Feira & Aula 07 & Business Intelligence, Random Forest e Ensemble Learning \\
\end{longtable}

\paragraph{OUTUBRO DE 2026}\label{outubro-de-2026}

\begin{longtable}[]{@{}
  >{\raggedright\arraybackslash}p{(\linewidth - 6\tabcolsep) * \real{0.0990}}
  >{\raggedright\arraybackslash}p{(\linewidth - 6\tabcolsep) * \real{0.1287}}
  >{\raggedright\arraybackslash}p{(\linewidth - 6\tabcolsep) * \real{0.0693}}
  >{\raggedright\arraybackslash}p{(\linewidth - 6\tabcolsep) * \real{0.7030}}@{}}
\toprule\noalign{}
\begin{minipage}[b]{\linewidth}\raggedright
Data
\end{minipage} & \begin{minipage}[b]{\linewidth}\raggedright
Dia da Semana
\end{minipage} & \begin{minipage}[b]{\linewidth}\raggedright
Aulas
\end{minipage} & \begin{minipage}[b]{\linewidth}\raggedright
Conteúdo
\end{minipage} \\
\midrule\noalign{}
\endhead
\bottomrule\noalign{}
\endlastfoot
06/10/2026 & Terça-Feira & Aula 08 & Big Data e Support Vector Machines \\
13/10/2026 & Terça-Feira & Aula 09 & Data-Driven Decision Making e Redução de Dimensionalidade \\
20/10/2026 & Terça-Feira & Aula 10 & Modelos de gestão de riscos e técnicas de Agrupamento \\
27/10/2026 & Terça-Feira & Aula 11 & Inteligência Artificial, análise preditiva e Seleção de Características \\
\end{longtable}

\paragraph{NOVEMBRO DE 2026}\label{novembro-de-2026}

\begin{longtable}[]{@{}
  >{\raggedright\arraybackslash}p{(\linewidth - 6\tabcolsep) * \real{0.0820}}
  >{\raggedright\arraybackslash}p{(\linewidth - 6\tabcolsep) * \real{0.1066}}
  >{\raggedright\arraybackslash}p{(\linewidth - 6\tabcolsep) * \real{0.0574}}
  >{\raggedright\arraybackslash}p{(\linewidth - 6\tabcolsep) * \real{0.7541}}@{}}
\toprule\noalign{}
\begin{minipage}[b]{\linewidth}\raggedright
Data
\end{minipage} & \begin{minipage}[b]{\linewidth}\raggedright
Dia da Semana
\end{minipage} & \begin{minipage}[b]{\linewidth}\raggedright
Aulas
\end{minipage} & \begin{minipage}[b]{\linewidth}\raggedright
Conteúdo
\end{minipage} \\
\midrule\noalign{}
\endhead
\bottomrule\noalign{}
\endlastfoot
03/11/2026 & Terça-Feira & Aula 12 & Monitoramento contínuo, melhoria de processos e ciclo de vida de modelos de Machine Learning \\
10/11/2026 & Terça-Feira & NP2 & PROVA \\
17/11/2026 & Terça-Feira & SUB & PROVA SUBSTITUTIVA \\
24/11/2026 & Terça-Feira & & N/A \\
\end{longtable}

\paragraph{DEZEMBRO DE 2026}\label{dezembro-de-2026}

\begin{longtable}[]{@{}llll@{}}
\toprule\noalign{}
Data & Dia da Semana & Aulas & Conteúdo \\
\midrule\noalign{}
\endhead
\bottomrule\noalign{}
\endlastfoot
01/12/2026 & Terça-Feira & & N/A \\
08/12/2026 & Terça-Feira & EXAME & PROVA \\
\end{longtable}

\chapter*{Apresentação da Disciplina}\label{apresentauxe7uxe3o-da-disciplina}
\addcontentsline{toc}{chapter}{Apresentação da Disciplina}

\chapter{- Dados com Ativo Estratégico: Introdução ao Machine Learning}\label{dados-com-ativo-estratuxe9gico-introduuxe7uxe3o-ao-machine-learning}

\section{- Exercícios - Dados com Atvo Estratégico: Introdução ao Machine Learning}\label{exercuxedcios---dados-com-atvo-estratuxe9gico-introduuxe7uxe3o-ao-machine-learning}

\chapter{- Governança de Dados e Ciclo de Vida de Machine Learning}\label{governanuxe7a-de-dados-e-ciclo-de-vida-de-machine-learning}

\section{- Governança de Dados e Ciclo de Vida de Machine Learning}\label{governanuxe7a-de-dados-e-ciclo-de-vida-de-machine-learning-1}

\chapter{Coleta de Dados e Regressão Linear}\label{coleta-de-dados-e-regressuxe3o-linear}

\pandocbounded{\includegraphics[keepaspectratio]{2026-AprendizadoProfundo_files/media/03-aula03/aula_regressao_linear/Slide1.JPG}} \citep[Capítulo 14][p.]{Grus2015DataScience}

\begin{figure}
\centering
\pandocbounded{\includegraphics[keepaspectratio]{2026-AprendizadoProfundo_files/media/03-aula03/aula_regressao_linear/Slide2.JPG}}
\caption{Aula 03 - Coleta de Dados e Regressão Linear}
\end{figure}

\begin{figure}
\centering
\pandocbounded{\includegraphics[keepaspectratio]{2026-AprendizadoProfundo_files/media/03-aula03/aula_regressao_linear/Slide3.JPG}}
\caption{Aula 03 - Coleta de Dados e Regressão Linear}
\end{figure}

\begin{figure}
\centering
\pandocbounded{\includegraphics[keepaspectratio]{2026-AprendizadoProfundo_files/media/03-aula03/aula_regressao_linear/Slide4.JPG}}
\caption{Aula 03 - Coleta de Dados e Regressão Linear}
\end{figure}

\begin{figure}
\centering
\pandocbounded{\includegraphics[keepaspectratio]{2026-AprendizadoProfundo_files/media/03-aula03/aula_regressao_linear/Slide5.JPG}}
\caption{Aula 03 - Coleta de Dados e Regressão Linear}
\end{figure}

\begin{figure}
\centering
\pandocbounded{\includegraphics[keepaspectratio]{2026-AprendizadoProfundo_files/media/03-aula03/aula_regressao_linear/Slide6.JPG}}
\caption{Aula 03 - Coleta de Dados e Regressão Linear}
\end{figure}

\begin{figure}
\centering
\pandocbounded{\includegraphics[keepaspectratio]{2026-AprendizadoProfundo_files/media/03-aula03/aula_regressao_linear/Slide7.JPG}}
\caption{Aula 03 - Coleta de Dados e Regressão Linear}
\end{figure}

\begin{figure}
\centering
\pandocbounded{\includegraphics[keepaspectratio]{2026-AprendizadoProfundo_files/media/03-aula03/aula_regressao_linear/Slide8.JPG}}
\caption{Aula 03 - Coleta de Dados e Regressão Linear}
\end{figure}

\begin{figure}
\centering
\pandocbounded{\includegraphics[keepaspectratio]{2026-AprendizadoProfundo_files/media/03-aula03/aula_regressao_linear/Slide9.JPG}}
\caption{Aula 03 - Coleta de Dados e Regressão Linear}
\end{figure}

\begin{figure}
\centering
\pandocbounded{\includegraphics[keepaspectratio]{2026-AprendizadoProfundo_files/media/03-aula03/aula_regressao_linear/Slide10.JPG}}
\caption{Aula 03 - Coleta de Dados e Regressão Linear}
\end{figure}

\begin{figure}
\centering
\pandocbounded{\includegraphics[keepaspectratio]{2026-AprendizadoProfundo_files/media/03-aula03/aula_regressao_linear/Slide11.JPG}}
\caption{Aula 03 - Coleta de Dados e Regressão Linear}
\end{figure}

\begin{figure}
\centering
\pandocbounded{\includegraphics[keepaspectratio]{2026-AprendizadoProfundo_files/media/03-aula03/aula_regressao_linear/Slide12.JPG}}
\caption{Aula 03 - Coleta de Dados e Regressão Linear}
\end{figure}

\begin{figure}
\centering
\pandocbounded{\includegraphics[keepaspectratio]{2026-AprendizadoProfundo_files/media/03-aula03/aula_regressao_linear/Slide13.JPG}}
\caption{Aula 03 - Coleta de Dados e Regressão Linear}
\end{figure}

\begin{figure}
\centering
\pandocbounded{\includegraphics[keepaspectratio]{2026-AprendizadoProfundo_files/media/03-aula03/aula_regressao_linear/Slide14.JPG}}
\caption{Aula 03 - Coleta de Dados e Regressão Linear}
\end{figure}

\begin{figure}
\centering
\pandocbounded{\includegraphics[keepaspectratio]{2026-AprendizadoProfundo_files/media/03-aula03/aula_regressao_linear/Slide15.JPG}}
\caption{Aula 03 - Coleta de Dados e Regressão Linear}
\end{figure}

\pandocbounded{\includegraphics[keepaspectratio]{2026-AprendizadoProfundo_files/2026-AprendizadoProfundo_files/figure-latex/linear-regression-labels-1.pdf}}

\begin{figure}
\centering
\pandocbounded{\includegraphics[keepaspectratio]{2026-AprendizadoProfundo_files/media/03-aula03/aula_regressao_linear/Slide16.JPG}}
\caption{Aula 03 - Coleta de Dados e Regressão Linear}
\end{figure}

\begin{figure}
\centering
\pandocbounded{\includegraphics[keepaspectratio]{2026-AprendizadoProfundo_files/media/03-aula03/aula_regressao_linear/Slide17.JPG}}
\caption{Aula 03 - Coleta de Dados e Regressão Linear}
\end{figure}

00000000000000000000000000000000000000000000000000000000000000000000000000000000

00000000000000000000000000000000000000000000000000000000000000000000000000000000

\begin{figure}
\centering
\pandocbounded{\includegraphics[keepaspectratio]{2026-AprendizadoProfundo_files/media/03-aula03/aula_regressao_linear/Slide18.JPG}}
\caption{Aula 03 - Coleta de Dados e Regressão Linear}
\end{figure}

00000000000000000000000000000000000000000000000000000000000000000000000000000000

00000000000000000000000000000000000000000000000000000000000000000000000000000000

\begin{figure}
\centering
\pandocbounded{\includegraphics[keepaspectratio]{2026-AprendizadoProfundo_files/media/03-aula03/aula_regressao_linear/Slide19.JPG}}
\caption{Aula 03 - Coleta de Dados e Regressão Linear}
\end{figure}

\begin{figure}
\centering
\pandocbounded{\includegraphics[keepaspectratio]{2026-AprendizadoProfundo_files/media/03-aula03/aula_regressao_linear/Slide20.JPG}}
\caption{Aula 03 - Coleta de Dados e Regressão Linear}
\end{figure}

\begin{figure}
\centering
\pandocbounded{\includegraphics[keepaspectratio]{2026-AprendizadoProfundo_files/media/03-aula03/aula_regressao_linear/Slide21.JPG}}
\caption{Aula 03 - Coleta de Dados e Regressão Linear}
\end{figure}

\begin{figure}
\centering
\pandocbounded{\includegraphics[keepaspectratio]{2026-AprendizadoProfundo_files/media/03-aula03/aula_regressao_linear/Slide22.JPG}}
\caption{Aula 03 - Coleta de Dados e Regressão Linear}
\end{figure}

\begin{figure}
\centering
\pandocbounded{\includegraphics[keepaspectratio]{2026-AprendizadoProfundo_files/media/03-aula03/aula_regressao_linear/Slide23.JPG}}
\caption{Aula 03 - Coleta de Dados e Regressão Linear}
\end{figure}

\begin{figure}
\centering
\pandocbounded{\includegraphics[keepaspectratio]{2026-AprendizadoProfundo_files/media/03-aula03/aula_regressao_linear/Slide24.JPG}}
\caption{Aula 03 - Coleta de Dados e Regressão Linear}
\end{figure}

\begin{Shaded}
\begin{Highlighting}[]
\ImportTok{import}\NormalTok{ numpy }\ImportTok{as}\NormalTok{ np}
\ImportTok{from}\NormalTok{ sklearn.linear\_model }\ImportTok{import}\NormalTok{ LinearRegression}

\NormalTok{X }\OperatorTok{=}\NormalTok{ np.array([[}\DecValTok{40}\NormalTok{], [}\DecValTok{50}\NormalTok{], [}\DecValTok{60}\NormalTok{], [}\DecValTok{70}\NormalTok{], [}\DecValTok{80}\NormalTok{]])}
\NormalTok{y }\OperatorTok{=}\NormalTok{ np.array([}\DecValTok{220}\NormalTok{, }\DecValTok{260}\NormalTok{, }\DecValTok{310}\NormalTok{, }\DecValTok{350}\NormalTok{, }\DecValTok{410}\NormalTok{])}

\CommentTok{\# Crie o seu modelo}
\NormalTok{modelo }\OperatorTok{=}\NormalTok{ LinearRegression()}

\CommentTok{\# Treine o seu modelo (ajuste sua reta aos pontos)}
\NormalTok{modelo.fit(X, y)}
\end{Highlighting}
\end{Shaded}

\begin{verbatim}
## LinearRegression()
\end{verbatim}

\begin{Shaded}
\begin{Highlighting}[]
\CommentTok{\# }
\BuiltInTok{print}\NormalTok{(}\StringTok{"O valor de W é: "}\NormalTok{,modelo.coef\_)}
\end{Highlighting}
\end{Shaded}

\begin{verbatim}
## O valor de W é:  [4.7]
\end{verbatim}

\begin{Shaded}
\begin{Highlighting}[]
\BuiltInTok{print}\NormalTok{(}\SpecialStringTok{f"O valor de b é }\SpecialCharTok{\{}\NormalTok{modelo}\SpecialCharTok{.}\NormalTok{intercept\_}\SpecialCharTok{\}}\SpecialStringTok{"}\NormalTok{)}
\end{Highlighting}
\end{Shaded}

\begin{verbatim}
## O valor de b é 28.0
\end{verbatim}

\begin{Shaded}
\begin{Highlighting}[]
\CommentTok{\# calcule: para um Apartamento de 65m\^{}2 qual é o valor previsto em mil reais ?}

\BuiltInTok{print}\NormalTok{(}\SpecialStringTok{f"Qunado X vale 65, a saida vale: }\SpecialCharTok{\{}\NormalTok{modelo}\SpecialCharTok{.}\NormalTok{predict([[}\DecValTok{65}\NormalTok{]])}\SpecialCharTok{\}}\SpecialStringTok{"}\NormalTok{)}
\end{Highlighting}
\end{Shaded}

\begin{verbatim}
## Qunado X vale 65, a saida vale: [333.5]
\end{verbatim}

\begin{figure}
\centering
\pandocbounded{\includegraphics[keepaspectratio]{2026-AprendizadoProfundo_files/media/03-aula03/aula_regressao_linear/Slide25.JPG}}
\caption{Aula 03 - Coleta de Dados e Regressão Linear}
\end{figure}

\begin{figure}
\centering
\pandocbounded{\includegraphics[keepaspectratio]{2026-AprendizadoProfundo_files/media/03-aula03/aula_regressao_linear/Slide26.JPG}}
\caption{Aula 03 - Coleta de Dados e Regressão Linear}
\end{figure}

\begin{figure}
\centering
\pandocbounded{\includegraphics[keepaspectratio]{2026-AprendizadoProfundo_files/media/03-aula03/aula_regressao_linear/Slide27.JPG}}
\caption{Aula 03 - Coleta de Dados e Regressão Linear}
\end{figure}

\citep{Grus2015DataScience}

\section{- Exercícios - Coleta de Dados e Regressão Linear}\label{exercuxedcios---coleta-de-dados-e-regressuxe3o-linear}

\subsection{Exemplo 01 -}\label{exemplo-01--}

Uma pesquisa foi realizada com o objetivo de verificar se existe
associação entre a falta de sono e a capacidade de as pessoas resolverem
problemas simples. Foram testadas 10 pessoas, mantendo-se sem dormir por
um determinado número de horas. Após cada um destes períodos, cada
pessoa teve de resolver um teste com adições simples, anotando-se então
os erros cometidos. Os dados resultantes são os seguintes:

\begin{longtable}[]{@{}
  >{\centering\arraybackslash}p{(\linewidth - 4\tabcolsep) * \real{0.3194}}
  >{\centering\arraybackslash}p{(\linewidth - 4\tabcolsep) * \real{0.3194}}
  >{\centering\arraybackslash}p{(\linewidth - 4\tabcolsep) * \real{0.3194}}@{}}
\toprule\noalign{}
\begin{minipage}[b]{\linewidth}\centering
Pessoa Observada (i)
\end{minipage} & \begin{minipage}[b]{\linewidth}\centering
\textbf{Número de Horas
sem Dormir ( X )}
\end{minipage} & \begin{minipage}[b]{\linewidth}\centering
\textbf{Número de Erros (
Y )}
\end{minipage} \\
\midrule\noalign{}
\endhead
\bottomrule\noalign{}
\endlastfoot
1 & 8 & 6 \\
2 & 8 & 8 \\
3 & 12 & 6 \\
4 & 12 & 10 \\
5 & 16 & 8 \\
6 & 16 & 14 \\
7 & 20 & 12 \\
8 & 20 & 14 \\
9 & 24 & 12 \\
10 & 24 & 16 \\
\end{longtable}

\begin{itemize}
\tightlist
\item
  \begin{enumerate}
  \def\labelenumi{\alph{enumi})}
  \tightlist
  \item
    Calcule o \textbf{coeficiente de correlação linear} de Pearson,
    interprete e teste sua significância a \textbf{5\%}.
  \end{enumerate}
\item
  \begin{enumerate}
  \def\labelenumi{\alph{enumi})}
  \setcounter{enumi}{1}
  \tightlist
  \item
    Ajuste, \textbf{teste e interprete o modelo de regressão linear}
    simples.
  \end{enumerate}
\end{itemize}

No python Comece importando os cabeçalhos: - importe a biblioteca
\textbf{numpy} para criar vetores de valores - importe a biblioteca \textbf{stats}
para usar as funções para \texttt{linregress} para calcular a regressão
linear - importe a biblioteca \textbf{stats} para usar as funções para
\texttt{pearsonr} para calcular o coeficiente de correlação \texttt{r}

\begin{Shaded}
\begin{Highlighting}[]
\ImportTok{import}\NormalTok{ numpy }\ImportTok{as}\NormalTok{ np}
\ImportTok{from}\NormalTok{ scipy }\ImportTok{import}\NormalTok{ stats}
\ImportTok{from}\NormalTok{ yt\_dlp.extractor }\ImportTok{import}\NormalTok{ agora}
\end{Highlighting}
\end{Shaded}

Agora precisamos criar uma função para calcular a \textbf{regressão linear}

\begin{Shaded}
\begin{Highlighting}[]
\KeywordTok{def}\NormalTok{ calcular\_regressao\_linear(x: np.ndarray, y: np.ndarray) }\OperatorTok{{-}\textgreater{}} \BuiltInTok{tuple}\NormalTok{[}\BuiltInTok{float}\NormalTok{, }\BuiltInTok{float}\NormalTok{, }\BuiltInTok{float}\NormalTok{]:}
    \CommentTok{"""Calcula os parâmetros e o teste de significância da regressão linear simples."""}
    \BuiltInTok{print}\NormalTok{(}\StringTok{"{-}{-}{-} b) Modelo de Regressão Linear Simples {-}{-}{-}"}\NormalTok{)}
    \CommentTok{\# Ajustando a reta de regressão e obtendo estatísticas completas}
\NormalTok{    inclinacao\_reta, interceptacao, correlacao\_r, valor\_p, erro\_padrao }\OperatorTok{=}\NormalTok{ stats.linregress(x, y)}
    \BuiltInTok{print}\NormalTok{(}\SpecialStringTok{f"Coeficiente Angular ( m ): }\SpecialCharTok{\{}\NormalTok{inclinacao\_reta}\SpecialCharTok{:.4f\}}\SpecialStringTok{"}\NormalTok{)}
    \BuiltInTok{print}\NormalTok{(}\SpecialStringTok{f"Intercepto ( n ): }\SpecialCharTok{\{}\NormalTok{interceptacao}\SpecialCharTok{:.4f\}}\SpecialStringTok{"}\NormalTok{)}
    \BuiltInTok{print}\NormalTok{(}\SpecialStringTok{f"Equação ajustada: Y\_hat = }\SpecialCharTok{\{}\NormalTok{interceptacao}\SpecialCharTok{:.4f\}}\SpecialStringTok{ + }\SpecialCharTok{\{}\NormalTok{inclinacao\_reta}\SpecialCharTok{:.4f\}}\SpecialStringTok{ * X"}\NormalTok{)}
    \BuiltInTok{print}\NormalTok{(}\SpecialStringTok{f"P{-}valor do coeficiente angular (teste do modelo): }\SpecialCharTok{\{}\NormalTok{valor\_p}\SpecialCharTok{:.5f\}}\SpecialStringTok{"}\NormalTok{)}

    \ControlFlowTok{if}\NormalTok{ ( valor\_p }\OperatorTok{\textless{}} \FloatTok{0.05}\NormalTok{):}
        \BuiltInTok{print}\NormalTok{(}\StringTok{"Conclusão: O modelo de regressão linear é estatisticamente significativo a 5\%."}\NormalTok{)}
    \ControlFlowTok{else}\NormalTok{:}
        \BuiltInTok{print}\NormalTok{(}\StringTok{"Conclusão: O modelo não é estatisticamente significativo."}\NormalTok{)}

    \ControlFlowTok{return}\NormalTok{ interceptacao, inclinacao\_reta, valor\_p}
\end{Highlighting}
\end{Shaded}

Agora precisamos criar uma função para calcular o \textbf{coeficiente de
correlação} \texttt{r}

\begin{Shaded}
\begin{Highlighting}[]
\KeywordTok{def}\NormalTok{ calcular\_correlacao\_pearson(v\_independentes\_x: np.ndarray, v\_dependentes\_y: np.ndarray) }\OperatorTok{{-}\textgreater{}} \BuiltInTok{tuple}\NormalTok{[}\BuiltInTok{float}\NormalTok{, }\BuiltInTok{float}\NormalTok{]:}
    \CommentTok{"""Calcula o coeficiente de correlação de Pearson e o p{-}valor correspondente."""}
    \BuiltInTok{print}\NormalTok{(}\StringTok{"{-}{-}{-} a) Coeficiente de Correlação de Pearson e Teste {-}{-}{-}"}\NormalTok{)}
\NormalTok{    coeficiente\_de\_correlacao\_r, valor\_de\_p }\OperatorTok{=}\NormalTok{ stats.pearsonr(v\_independentes\_x, v\_dependentes\_y)}

    \BuiltInTok{print}\NormalTok{(}\SpecialStringTok{f"Coeficiente de Correlação (r): }\SpecialCharTok{\{}\NormalTok{coeficiente\_de\_correlacao\_r}\SpecialCharTok{:.4f\}}\SpecialStringTok{"}\NormalTok{)}
    \BuiltInTok{print}\NormalTok{(}\SpecialStringTok{f"P{-}valor da correlação: }\SpecialCharTok{\{}\NormalTok{valor\_de\_p}\SpecialCharTok{:.5f\}}\SpecialStringTok{"}\NormalTok{)}

    \ControlFlowTok{if}\NormalTok{ ( valor\_de\_p }\OperatorTok{\textless{}} \FloatTok{0.05}\NormalTok{):}
        \BuiltInTok{print}\NormalTok{(}\StringTok{"Conclusão: Rejeitamos a Hipótese. A correlação é estatisticamente relevante a 5\%.}\CharTok{\textbackslash{}n}\StringTok{"}\NormalTok{)}
    \ControlFlowTok{else}\NormalTok{:}
        \BuiltInTok{print}\NormalTok{(}\StringTok{"Conclusão: Não rejeitamos a Hipótese.}\CharTok{\textbackslash{}n}\StringTok{"}\NormalTok{)}
    \ControlFlowTok{return}\NormalTok{ coeficiente\_de\_correlacao\_r, valor\_de\_p}
\end{Highlighting}
\end{Shaded}

Agora vamos definir a ordem de execução das funções:

\begin{Shaded}
\begin{Highlighting}[]
\KeywordTok{def}\NormalTok{ funcao\_principal() }\OperatorTok{{-}\textgreater{}} \VariableTok{None}\NormalTok{:}
    \CommentTok{\# Carregar os valores de X (que é a variável independente) no vetor variaveis\_independentes\_x}
    \CommentTok{\# X: Número de horas sem dormir}
\NormalTok{    variaveis\_independentes\_x }\OperatorTok{=}\NormalTok{ np.array([}\DecValTok{8}\NormalTok{, }\DecValTok{8}\NormalTok{, }\DecValTok{12}\NormalTok{, }\DecValTok{12}\NormalTok{, }\DecValTok{16}\NormalTok{, }\DecValTok{16}\NormalTok{, }\DecValTok{20}\NormalTok{, }\DecValTok{20}\NormalTok{, }\DecValTok{24}\NormalTok{, }\DecValTok{24}\NormalTok{])}

    \CommentTok{\# \# Carregar os valores de Y (que é a variável independente) no vetor variaveis\_dependentes\_y}
    \CommentTok{\# Y: Número de erros}
\NormalTok{    variaveis\_dependentes\_y }\OperatorTok{=}\NormalTok{ np.array([}\DecValTok{6}\NormalTok{, }\DecValTok{8}\NormalTok{, }\DecValTok{6}\NormalTok{, }\DecValTok{10}\NormalTok{, }\DecValTok{8}\NormalTok{, }\DecValTok{14}\NormalTok{, }\DecValTok{12}\NormalTok{, }\DecValTok{14}\NormalTok{, }\DecValTok{12}\NormalTok{, }\DecValTok{16}\NormalTok{])}

    \CommentTok{\# Executando a correlação}
\NormalTok{    calcular\_correlacao\_pearson(variaveis\_independentes\_x, variaveis\_dependentes\_y)}

    \CommentTok{\# Executando a regressão linear}
\NormalTok{    calcular\_regressao\_linear(variaveis\_independentes\_x, variaveis\_dependentes\_y)}
\end{Highlighting}
\end{Shaded}

Finalmente, vamos iniciar o \textbf{fluxo principal} do programa, chamado a
\textbf{funcao\_principal()} que chamará todas as outras na ordem

\begin{Shaded}
\begin{Highlighting}[]
\ControlFlowTok{if} \VariableTok{\_\_name\_\_} \OperatorTok{==} \StringTok{\textquotesingle{}\_\_main\_\_\textquotesingle{}}\NormalTok{:}
\NormalTok{    funcao\_principal()}
\end{Highlighting}
\end{Shaded}

\begin{verbatim}
## --- a) Coeficiente de Correlação de Pearson e Teste ---
## Coeficiente de Correlação (r): 0.8015
## P-valor da correlação: 0.00531
## Conclusão: Rejeitamos a Hipótese. A correlação é estatisticamente relevante a 5%.
## 
## --- b) Modelo de Regressão Linear Simples ---
## Coeficiente Angular ( m ): 0.4750
## Intercepto ( n ): 3.0000
## Equação ajustada: Y_hat = 3.0000 + 0.4750 * X
## P-valor do coeficiente angular (teste do modelo): 0.00531
## Conclusão: O modelo de regressão linear é estatisticamente significativo a 5%.
\end{verbatim}

Como ficou a equação descoberta pela \textbf{regressão linear simples} ?

Uma equação de reta segue o formato \(y = mx + n\).

\[
y = \underbrace{m}*{*\text{coeficiente angular}} x + \underbrace{n}{\text{coeficiente linear}}
\]

No exercício: - o valor de \texttt{m} é dado pelo valor do
\texttt{coeficiente\ angular} ; - o valor de \texttt{n} é dado pelo valor do
\texttt{intercepto} ;

Então a equação ficou:

\[
y = 0,4750(x) + 3
\]

\subsubsection{Gráfico do Exemplo 00}\label{gruxe1fico-do-exemplo-00}

\begin{Shaded}
\begin{Highlighting}[]
\ImportTok{import}\NormalTok{ numpy }\ImportTok{as}\NormalTok{ np}
\ImportTok{import}\NormalTok{ matplotlib.pyplot }\ImportTok{as}\NormalTok{ plt}

\CommentTok{\# Dados do problema}
\NormalTok{x }\OperatorTok{=}\NormalTok{ np.array([}\DecValTok{8}\NormalTok{, }\DecValTok{8}\NormalTok{, }\DecValTok{12}\NormalTok{, }\DecValTok{12}\NormalTok{, }\DecValTok{16}\NormalTok{, }\DecValTok{16}\NormalTok{, }\DecValTok{20}\NormalTok{, }\DecValTok{20}\NormalTok{, }\DecValTok{24}\NormalTok{, }\DecValTok{24}\NormalTok{])}
\NormalTok{y }\OperatorTok{=}\NormalTok{ np.array([}\DecValTok{6}\NormalTok{, }\DecValTok{8}\NormalTok{, }\DecValTok{6}\NormalTok{, }\DecValTok{10}\NormalTok{, }\DecValTok{8}\NormalTok{, }\DecValTok{14}\NormalTok{, }\DecValTok{12}\NormalTok{, }\DecValTok{14}\NormalTok{, }\DecValTok{12}\NormalTok{, }\DecValTok{16}\NormalTok{])}

\CommentTok{\# Configuração do gráfico}
\NormalTok{plt.figure(figsize}\OperatorTok{=}\NormalTok{(}\DecValTok{10}\NormalTok{, }\DecValTok{7}\NormalTok{))}

\CommentTok{\# Definindo os limites dos eixos começando em 0}
\NormalTok{plt.xlim(}\OperatorTok{{-}}\DecValTok{8}\NormalTok{, }\BuiltInTok{max}\NormalTok{(x) }\OperatorTok{+} \DecValTok{2}\NormalTok{)}
\end{Highlighting}
\end{Shaded}

\begin{verbatim}
## (-8.0, 26.0)
\end{verbatim}

\begin{Shaded}
\begin{Highlighting}[]
\NormalTok{plt.ylim(}\OperatorTok{{-}}\DecValTok{8}\NormalTok{, }\BuiltInTok{max}\NormalTok{(y) }\OperatorTok{+} \DecValTok{2}\NormalTok{)}
\end{Highlighting}
\end{Shaded}

\begin{verbatim}
## (-8.0, 18.0)
\end{verbatim}

\begin{Shaded}
\begin{Highlighting}[]
\CommentTok{\# Configurando a escala unitária (passo de 1 em 1) para os eixos X e Y}
\NormalTok{plt.xticks(np.arange(}\DecValTok{0}\NormalTok{, }\BuiltInTok{max}\NormalTok{(x) }\OperatorTok{+} \DecValTok{3}\NormalTok{, }\DecValTok{1}\NormalTok{))}
\end{Highlighting}
\end{Shaded}

\begin{verbatim}
## ([<matplotlib.axis.XTick object at 0x000001CC61887530>, <matplotlib.axis.XTick object at 0x000001CC61AAD1F0>, <matplotlib.axis.XTick object at 0x000001CC61A03C50>, <matplotlib.axis.XTick object at 0x000001CC61B2CCB0>, <matplotlib.axis.XTick object at 0x000001CC61A8FBC0>, <matplotlib.axis.XTick object at 0x000001CC61A8EAB0>, <matplotlib.axis.XTick object at 0x000001CC61B2D5E0>, <matplotlib.axis.XTick object at 0x000001CC61B2E0F0>, <matplotlib.axis.XTick object at 0x000001CC61B2ED20>, <matplotlib.axis.XTick object at 0x000001CC61B2E360>, <matplotlib.axis.XTick object at 0x000001CC61B30FB0>, <matplotlib.axis.XTick object at 0x000001CC61B32DB0>, <matplotlib.axis.XTick object at 0x000001CC61B33770>, <matplotlib.axis.XTick object at 0x000001CC61B38230>, <matplotlib.axis.XTick object at 0x000001CC61B33110>, <matplotlib.axis.XTick object at 0x000001CC61B38DD0>, <matplotlib.axis.XTick object at 0x000001CC61B398B0>, <matplotlib.axis.XTick object at 0x000001CC61B3A3F0>, <matplotlib.axis.XTick object at 0x000001CC61B3A840>, <matplotlib.axis.XTick object at 0x000001CC61B38A10>, <matplotlib.axis.XTick object at 0x000001CC61B3FB90>, <matplotlib.axis.XTick object at 0x000001CC619CA8D0>, <matplotlib.axis.XTick object at 0x000001CC619CB3B0>, <matplotlib.axis.XTick object at 0x000001CC619CBC20>, <matplotlib.axis.XTick object at 0x000001CC61B3F560>, <matplotlib.axis.XTick object at 0x000001CC61C2AC30>, <matplotlib.axis.XTick object at 0x000001CC61C2B800>], [Text(0, 0, '0'), Text(1, 0, '1'), Text(2, 0, '2'), Text(3, 0, '3'), Text(4, 0, '4'), Text(5, 0, '5'), Text(6, 0, '6'), Text(7, 0, '7'), Text(8, 0, '8'), Text(9, 0, '9'), Text(10, 0, '10'), Text(11, 0, '11'), Text(12, 0, '12'), Text(13, 0, '13'), Text(14, 0, '14'), Text(15, 0, '15'), Text(16, 0, '16'), Text(17, 0, '17'), Text(18, 0, '18'), Text(19, 0, '19'), Text(20, 0, '20'), Text(21, 0, '21'), Text(22, 0, '22'), Text(23, 0, '23'), Text(24, 0, '24'), Text(25, 0, '25'), Text(26, 0, '26')])
\end{verbatim}

\begin{Shaded}
\begin{Highlighting}[]
\NormalTok{plt.yticks(np.arange(}\DecValTok{0}\NormalTok{, }\BuiltInTok{max}\NormalTok{(y) }\OperatorTok{+} \DecValTok{3}\NormalTok{, }\DecValTok{1}\NormalTok{))}
\end{Highlighting}
\end{Shaded}

\begin{verbatim}
## ([<matplotlib.axis.YTick object at 0x000001CC61AAC620>, <matplotlib.axis.YTick object at 0x000001CC61AAC500>, <matplotlib.axis.YTick object at 0x000001CC61C28140>, <matplotlib.axis.YTick object at 0x000001CC61C2D160>, <matplotlib.axis.YTick object at 0x000001CC61C2E2A0>, <matplotlib.axis.YTick object at 0x000001CC61C2E990>, <matplotlib.axis.YTick object at 0x000001CC61C2F830>, <matplotlib.axis.YTick object at 0x000001CC61C2E900>, <matplotlib.axis.YTick object at 0x000001CC61C2FE30>, <matplotlib.axis.YTick object at 0x000001CC61C30C50>, <matplotlib.axis.YTick object at 0x000001CC61C316D0>, <matplotlib.axis.YTick object at 0x000001CC61C32210>, <matplotlib.axis.YTick object at 0x000001CC61C310D0>, <matplotlib.axis.YTick object at 0x000001CC61C32C00>, <matplotlib.axis.YTick object at 0x000001CC61C33770>, <matplotlib.axis.YTick object at 0x000001CC61C34290>, <matplotlib.axis.YTick object at 0x000001CC61C34CE0>, <matplotlib.axis.YTick object at 0x000001CC61C32840>, <matplotlib.axis.YTick object at 0x000001CC61936060>], [Text(0, 0, '0'), Text(0, 1, '1'), Text(0, 2, '2'), Text(0, 3, '3'), Text(0, 4, '4'), Text(0, 5, '5'), Text(0, 6, '6'), Text(0, 7, '7'), Text(0, 8, '8'), Text(0, 9, '9'), Text(0, 10, '10'), Text(0, 11, '11'), Text(0, 12, '12'), Text(0, 13, '13'), Text(0, 14, '14'), Text(0, 15, '15'), Text(0, 16, '16'), Text(0, 17, '17'), Text(0, 18, '18')])
\end{verbatim}

\begin{Shaded}
\begin{Highlighting}[]
\CommentTok{\# Grade padrão em cinza claro para a escala geral}
\ControlFlowTok{for}\NormalTok{ xtick }\KeywordTok{in}\NormalTok{ np.arange(}\DecValTok{0}\NormalTok{, }\BuiltInTok{max}\NormalTok{(x) }\OperatorTok{+} \DecValTok{3}\NormalTok{, }\DecValTok{1}\NormalTok{):}
\NormalTok{    plt.axvline(xtick, color}\OperatorTok{=}\StringTok{\textquotesingle{}gray\textquotesingle{}}\NormalTok{, linewidth}\OperatorTok{=}\FloatTok{0.3}\NormalTok{, linestyle}\OperatorTok{=}\StringTok{\textquotesingle{}:\textquotesingle{}}\NormalTok{, alpha}\OperatorTok{=}\FloatTok{0.5}\NormalTok{, zorder}\OperatorTok{=}\DecValTok{1}\NormalTok{)}

\ControlFlowTok{for}\NormalTok{ ytick }\KeywordTok{in}\NormalTok{ np.arange(}\DecValTok{0}\NormalTok{, }\BuiltInTok{max}\NormalTok{(y) }\OperatorTok{+} \DecValTok{3}\NormalTok{, }\DecValTok{1}\NormalTok{):}
\NormalTok{    plt.axhline(ytick, color}\OperatorTok{=}\StringTok{\textquotesingle{}gray\textquotesingle{}}\NormalTok{, linewidth}\OperatorTok{=}\FloatTok{0.3}\NormalTok{, linestyle}\OperatorTok{=}\StringTok{\textquotesingle{}:\textquotesingle{}}\NormalTok{, alpha}\OperatorTok{=}\FloatTok{0.5}\NormalTok{, zorder}\OperatorTok{=}\DecValTok{1}\NormalTok{)}

\CommentTok{\# Projetando linhas vermelhas contínuas sobre os pares ordenados formados pelos vetores x e y}
\ControlFlowTok{for}\NormalTok{ xi, yi }\KeywordTok{in} \BuiltInTok{zip}\NormalTok{(x, y):}
\NormalTok{    plt.vlines(xi, }\DecValTok{0}\NormalTok{, yi, color}\OperatorTok{=}\StringTok{\textquotesingle{}green\textquotesingle{}}\NormalTok{, linewidth}\OperatorTok{=}\FloatTok{0.5}\NormalTok{, linestyle}\OperatorTok{=}\StringTok{\textquotesingle{}{-}{-}\textquotesingle{}}\NormalTok{, alpha}\OperatorTok{=}\FloatTok{0.5}\NormalTok{, zorder}\OperatorTok{=}\DecValTok{2}\NormalTok{)}
\NormalTok{    plt.hlines(yi, }\DecValTok{0}\NormalTok{, xi, color}\OperatorTok{=}\StringTok{\textquotesingle{}green\textquotesingle{}}\NormalTok{, linewidth}\OperatorTok{=}\FloatTok{0.5}\NormalTok{, linestyle}\OperatorTok{=}\StringTok{\textquotesingle{}{-}{-}\textquotesingle{}}\NormalTok{, alpha}\OperatorTok{=}\FloatTok{0.5}\NormalTok{, zorder}\OperatorTok{=}\DecValTok{2}\NormalTok{)}

\CommentTok{\# Destacando os eixos principais em (0,0)}
\NormalTok{plt.axhline(}\DecValTok{0}\NormalTok{, color}\OperatorTok{=}\StringTok{\textquotesingle{}black\textquotesingle{}}\NormalTok{, linewidth}\OperatorTok{=}\FloatTok{1.2}\NormalTok{, zorder}\OperatorTok{=}\DecValTok{3}\NormalTok{)}
\NormalTok{plt.axvline(}\DecValTok{0}\NormalTok{, color}\OperatorTok{=}\StringTok{\textquotesingle{}black\textquotesingle{}}\NormalTok{, linewidth}\OperatorTok{=}\FloatTok{1.2}\NormalTok{, zorder}\OperatorTok{=}\DecValTok{3}\NormalTok{)}

\CommentTok{\# Gráfico de dispersão (pontos observados)}
\NormalTok{plt.scatter(x, y, color}\OperatorTok{=}\StringTok{\textquotesingle{}blue\textquotesingle{}}\NormalTok{, label}\OperatorTok{=}\StringTok{\textquotesingle{}Dados Observados\textquotesingle{}}\NormalTok{, s}\OperatorTok{=}\DecValTok{60}\NormalTok{, zorder}\OperatorTok{=}\DecValTok{5}\NormalTok{)}

\CommentTok{\# Adicionando rótulos com os pares ordenados (X, Y) em cada ponto}
\ControlFlowTok{for}\NormalTok{ xi, yi }\KeywordTok{in} \BuiltInTok{zip}\NormalTok{(x, y):}
\NormalTok{    plt.annotate(}\SpecialStringTok{f"(}\SpecialCharTok{\{}\NormalTok{xi}\SpecialCharTok{\}}\SpecialStringTok{, }\SpecialCharTok{\{}\NormalTok{yi}\SpecialCharTok{\}}\SpecialStringTok{)"}\NormalTok{,}
\NormalTok{                 (xi, yi),}
\NormalTok{                 textcoords}\OperatorTok{=}\StringTok{"offset points"}\NormalTok{,}
\NormalTok{                 xytext}\OperatorTok{=}\NormalTok{(}\DecValTok{0}\NormalTok{, }\DecValTok{8}\NormalTok{),  }\CommentTok{\# Deslocamento vertical do texto em relação ao ponto}
\NormalTok{                 ha}\OperatorTok{=}\StringTok{\textquotesingle{}center\textquotesingle{}}\NormalTok{,}
\NormalTok{                 fontsize}\OperatorTok{=}\DecValTok{8}\NormalTok{,}
\NormalTok{                 fontweight}\OperatorTok{=}\StringTok{\textquotesingle{}bold\textquotesingle{}}\NormalTok{,}
\NormalTok{                 color}\OperatorTok{=}\StringTok{\textquotesingle{}darkblue\textquotesingle{}}\NormalTok{,}
\NormalTok{                 zorder}\OperatorTok{=}\DecValTok{6}\NormalTok{)}

\CommentTok{\# Definindo os pontos para a reta y = 0.4750 * x + 3}
\NormalTok{x\_linha }\OperatorTok{=}\NormalTok{ np.linspace(}\OperatorTok{{-}}\DecValTok{7}\NormalTok{, }\BuiltInTok{max}\NormalTok{(x), }\DecValTok{100}\NormalTok{)}
\NormalTok{y\_linha }\OperatorTok{=} \FloatTok{0.4750} \OperatorTok{*}\NormalTok{ x\_linha }\OperatorTok{+} \DecValTok{3}

\CommentTok{\# Plotando a reta de regressão}
\NormalTok{plt.plot(x\_linha, y\_linha, color}\OperatorTok{=}\StringTok{\textquotesingle{}red\textquotesingle{}}\NormalTok{, linestyle}\OperatorTok{=}\StringTok{\textquotesingle{}{-}\textquotesingle{}}\NormalTok{, linewidth}\OperatorTok{=}\DecValTok{2}\NormalTok{, label}\OperatorTok{=}\StringTok{\textquotesingle{}Reta: $Y = 0,4750X + 3$\textquotesingle{}}\NormalTok{, zorder}\OperatorTok{=}\DecValTok{4}\NormalTok{)}

\CommentTok{\# Rótulos e formatação}
\NormalTok{plt.title(}\StringTok{\textquotesingle{}Gráfico com Pares Ordenados (X, Y) Identificados nos Pontos\textquotesingle{}}\NormalTok{, fontsize}\OperatorTok{=}\DecValTok{12}\NormalTok{, fontweight}\OperatorTok{=}\StringTok{\textquotesingle{}bold\textquotesingle{}}\NormalTok{)}
\NormalTok{plt.xlabel(}\StringTok{\textquotesingle{}Número de Horas sem Dormir (X)\textquotesingle{}}\NormalTok{, fontsize}\OperatorTok{=}\DecValTok{10}\NormalTok{)}
\NormalTok{plt.ylabel(}\StringTok{\textquotesingle{}Número de Erros (Y)\textquotesingle{}}\NormalTok{, fontsize}\OperatorTok{=}\DecValTok{10}\NormalTok{)}
\NormalTok{plt.legend()}

\CommentTok{\# Exibir o gráfico}
\NormalTok{plt.show()}
\end{Highlighting}
\end{Shaded}

\pandocbounded{\includegraphics[keepaspectratio]{2026-AprendizadoProfundo_files/2026-AprendizadoProfundo_files/figure-latex/exemplo00-grafico-1.pdf}}

\subsection{Testes de Multipla-Escolha:}\label{testes-de-multipla-escolha}

\section{Questão 1}\label{questuxe3o-1}

\begin{longtable}[]{@{}
  >{\raggedright\arraybackslash}p{(\linewidth - 0\tabcolsep) * \real{1.0000}}@{}}
\toprule\noalign{}
\begin{minipage}[b]{\linewidth}\raggedright
\textbf{Questão 1}
\end{minipage} \\
\midrule\noalign{}
\endhead
\bottomrule\noalign{}
\endlastfoot
Na \textbf{Regressão Linear simples}, o modelo é representado por:
\(\hat{y}=w(x)+b\)

Assinale a alternativa correta sobre os \textbf{parâmetros} \(w\) e \(b\). \\
\textbf{A)} \(w\) representa o erro do modelo e \(b\) representa o valor real
observado. \\
\textbf{B)} \(w\) representa a inclinação da reta e \(b\) representa o
intercepto/deslocamento vertical. \\
\textbf{C)} \(w\) representa a variável dependente e \(b\) representa a
variável independente. \\
\textbf{D)} \(w\) representa o número de observações e \(b\) representa o
número de previsões. \\
\textbf{E)} \(w\) e \(b\) são sempre fornecidos previamente e não são
aprendidos pelo modelo. \\
\end{longtable}

\begin{center}\rule{0.5\linewidth}{0.5pt}\end{center}

\begin{longtable}[]{@{}
  >{\raggedright\arraybackslash}p{(\linewidth - 0\tabcolsep) * \real{1.0000}}@{}}
\toprule\noalign{}
\begin{minipage}[b]{\linewidth}\raggedright
\textbf{Questão 2}
\end{minipage} \\
\midrule\noalign{}
\endhead
\bottomrule\noalign{}
\endlastfoot
Na \textbf{programação tradicional}, conhecemos uma entrada \(\textbf{x}\) e
uma função \(\textbf{f(x)}\), e então calculamos a saída \(\textbf{y}\).

No \textbf{aprendizado de máquina supervisionado}, durante o treinamento,
a lógica é diferente. Assinale a alternativa que melhor representa
esse processo. \\
\textbf{A)} Dados + modelo conhecido \(\rightarrow\) resposta. \\
\textbf{B)} Modelo + resposta \(\rightarrow\) dados. \\
\textbf{C)} Dados + respostas conhecidas \(\rightarrow\) aprendizado de um
modelo. \\
\textbf{D)} Dados sem respostas \(\rightarrow\) resultado exato. \\
\textbf{E)} Apenas a variável \(x\) é suficiente para descobrir qualquer
modelo. \\
\end{longtable}

\begin{center}\rule{0.5\linewidth}{0.5pt}\end{center}

\section{Questão 3}\label{questuxe3o-3}

\begin{longtable}[]{@{}
  >{\raggedright\arraybackslash}p{(\linewidth - 0\tabcolsep) * \real{1.0000}}@{}}
\toprule\noalign{}
\begin{minipage}[b]{\linewidth}\raggedright
\textbf{Questão 3}
\end{minipage} \\
\midrule\noalign{}
\endhead
\bottomrule\noalign{}
\endlastfoot
Considere o modelo de \textbf{Regressão Linear}:

\(\hat{y}=4{,}7x+28\)

onde:

• \(x\) representa a área do imóvel em \(m^2\);

• \(\hat{y}\) representa o preço estimado em milhares de reais.

Para um apartamento com \(65\,m^2\):

\(\hat{y}=4{,}7(65)+28\)

\(\hat{y}=333{,}5\)

Analise as afirmações:

\textbf{I.} O preço estimado para um apartamento de \(65\,m^2\) é de
aproximadamente \textbf{R\$ 333,5 mil}.

\textbf{II.} O coeficiente \(4{,}7\) indica que, segundo o modelo, cada
aumento de \(1\,m^2\) está associado a um aumento estimado de \textbf{R\$ 4,7
mil}.

\textbf{III.} O valor \(28\) corresponde à inclinação da reta.

Assinale a alternativa \textbf{correta}. \\
\textbf{A)} Apenas I está correta. \\
\textbf{B)} Apenas II está correta. \\
\textbf{C)} Apenas I e II estão corretas. \\
\textbf{D)} Apenas II e III estão corretas. \\
\textbf{E)} I, II e III estão corretas. \\
\end{longtable}

\begin{center}\rule{0.5\linewidth}{0.5pt}\end{center}

\section{Questão 4}\label{questuxe3o-4}

\begin{longtable}[]{@{}
  >{\raggedright\arraybackslash}p{(\linewidth - 0\tabcolsep) * \real{1.0000}}@{}}
\toprule\noalign{}
\begin{minipage}[b]{\linewidth}\raggedright
\textbf{Questão 4}
\end{minipage} \\
\midrule\noalign{}
\endhead
\bottomrule\noalign{}
\endlastfoot
Considere os seguintes \textbf{pares ordenados} \((x,y)\), ou ainda
\textbf{(entrada, saída)}, observados:

\((40,220),\ (50,260),\ (60,310),\ (70,350),\ (80,410)\)

A partir desses dados, foi obtido o \textbf{modelo}:

\(\hat{y}=4{,}7x+28\)

Analise as afirmações:

\textbf{I.} A variável \(x\) representa a área do apartamento.

\textbf{II.} A variável \(\hat{y}\) representa o preço estimado pelo modelo.

\textbf{III.} Depois que o modelo foi encontrado, ele pode ser utilizado
para estimar o preço de um apartamento com área diferente das que
aparecem na tabela, como \(65\,m^2\).

Assinale a alternativa \textbf{correta}. \\
\textbf{A)} Apenas I está correta. \\
\textbf{B)} Apenas I e II estão corretas. \\
\textbf{C)} Apenas II e III estão corretas. \\
\textbf{D)} Apenas I e III estão corretas. \\
\textbf{E)} I, II e III estão corretas. \\
\end{longtable}

\begin{center}\rule{0.5\linewidth}{0.5pt}\end{center}

\chapter{Gabarito}\label{gabarito}

\begin{longtable}[]{@{}cc@{}}
\toprule\noalign{}
\textbf{Questão} & \textbf{Resposta correta} \\
\midrule\noalign{}
\endhead
\bottomrule\noalign{}
\endlastfoot
\textbf{1} & \textbf{B} \\
\textbf{2} & \textbf{C} \\
\textbf{3} & \textbf{C} \\
\textbf{4} & \textbf{E} \\
\end{longtable}

Também retirei da Questão 2 o comentário explicativo após o gabarito,
para manter todas as questões no mesmo padrão de prova.

\chapter{- Organização e Preparação dos Dados, Regressão Logística e KNN}\label{organizauxe7uxe3o-e-preparauxe7uxe3o-dos-dados-regressuxe3o-loguxedstica-e-knn}

\section{- Exercícios - Organização e Preparação dos Dados, Regressão Logística e KNN}\label{exercuxedcios---organizauxe7uxe3o-e-preparauxe7uxe3o-dos-dados-regressuxe3o-loguxedstica-e-knn}

\subsection{Organização e Preparação de Dados: População, Estrato , proporção e Amostra}\label{organizauxe7uxe3o-e-preparauxe7uxe3o-de-dados-populauxe7uxe3o-estrato-proporuxe7uxe3o-e-amostra}

A regra é:

\[
n_h = n\cdot\frac{N_h}{N}
\]

onde (N) é o tamanho da população, (N\_h) é o tamanho do estrato, (n) é o tamanho total da amostra e (n\_h) é a quantidade de elementos sorteados naquele estrato.

\subsubsection{1. Alunos de uma universidade}\label{alunos-de-uma-universidade}

\textbf{População:} 10.000 alunos.
\textbf{Amostra desejada:} 1.000 alunos.

\begin{longtable}[]{@{}
  >{\raggedright\arraybackslash}p{(\linewidth - 6\tabcolsep) * \real{0.2025}}
  >{\raggedleft\arraybackslash}p{(\linewidth - 6\tabcolsep) * \real{0.2911}}
  >{\raggedleft\arraybackslash}p{(\linewidth - 6\tabcolsep) * \real{0.1899}}
  >{\raggedleft\arraybackslash}p{(\linewidth - 6\tabcolsep) * \real{0.2911}}@{}}
\toprule\noalign{}
\begin{minipage}[b]{\linewidth}\raggedright
Estrato
\end{minipage} & \begin{minipage}[b]{\linewidth}\raggedleft
População no estrato
\end{minipage} & \begin{minipage}[b]{\linewidth}\raggedleft
Proporção
\end{minipage} & \begin{minipage}[b]{\linewidth}\raggedleft
Amostra proporcional
\end{minipage} \\
\midrule\noalign{}
\endhead
\bottomrule\noalign{}
\endlastfoot
Engenharia & 3.000 & 30\% & 300 \\
Administração & 2.500 & 25\% & \_\_\_ \\
Direito & 2.000 & 20\% & \_\_\_ \\
Saúde & 1.500 & 15\% & \_\_\_ \\
Outros cursos & 1.000 & 10\% & \_\_\_ \\
\textbf{Total} & \textbf{10.000} & \textbf{100\%} & \textbf{1.000} \\
\end{longtable}

Por exemplo:

\[
n_{\text{Engenharia}} = 1000\cdot\frac{3000}{10000} = 300
\]

\begin{center}\rule{0.5\linewidth}{0.5pt}\end{center}

\subsubsection{2. Funcionários de uma empresa}\label{funcionuxe1rios-de-uma-empresa}

\textbf{População:} 2.000 funcionários.
\textbf{Amostra:} 200 funcionários.

\begin{longtable}[]{@{}lrrr@{}}
\toprule\noalign{}
Estrato & População & Proporção & Amostra \\
\midrule\noalign{}
\endhead
\bottomrule\noalign{}
\endlastfoot
Produção & 900 & 45\% & 90 \\
Administrativo & 400 & 20\% & \_\_ \\
Comercial & 300 & 15\% & \_\_ \\
Tecnologia & 250 & 12,5\% & \_\_ \\
Diretoria/Gestão & 150 & 7,5\% & \_\_ \\
\textbf{Total} & \textbf{2.000} & \textbf{100\%} & \textbf{200} \\
\end{longtable}

Assim, embora a área de Produção tenha muito mais funcionários, ela não domina arbitrariamente a pesquisa: sua participação na amostra permanece igual à sua participação na população.

\begin{center}\rule{0.5\linewidth}{0.5pt}\end{center}

\subsubsection{3. Moradores de uma cidade}\label{moradores-de-uma-cidade}

\textbf{População:} 50.000 moradores.
\textbf{Amostra:} 1.000 moradores.

Estratificação por faixa etária:

\begin{longtable}[]{@{}lrrr@{}}
\toprule\noalign{}
Estrato & População & Proporção & Amostra \\
\midrule\noalign{}
\endhead
\bottomrule\noalign{}
\endlastfoot
0--17 anos & 10.000 & 20\% & 200 \\
18--29 anos & 8.000 & 16\% & \_\_\_ \\
30--44 anos & 12.000 & 24\% & \_\_\_ \\
45--59 anos & 11.000 & 22\% & \_\_\_ \\
60 anos ou mais & 9.000 & 18\% & \_\_\_ \\
\textbf{Total} & \textbf{50.000} & \textbf{100\%} & \textbf{1.000} \\
\end{longtable}

Para a faixa de 30 a 44 anos:

\[
n_h = 1000\cdot\frac{12000}{50000}= 240
\]

\begin{center}\rule{0.5\linewidth}{0.5pt}\end{center}

\subsubsection{4. Imagens do CIFAR-10}\label{imagens-do-cifar-10}

Aqui temos um exemplo particularmente simples porque o CIFAR-10 é \textbf{balanceado entre as classes}.

\textbf{População:} 60.000 imagens.
\textbf{Amostra:} 6.000 imagens.

Cada uma das 10 classes possui 6.000 imagens.

\begin{longtable}[]{@{}lrrr@{}}
\toprule\noalign{}
Estrato --- classe & População & Proporção & Amostra \\
\midrule\noalign{}
\endhead
\bottomrule\noalign{}
\endlastfoot
airplane & 6.000 & 10\% & 600 \\
automobile & 6.000 & 10\% & \_\_\_ \\
bird & 6.000 & 10\% & \_\_\_ \\
cat & 6.000 & 10\% & \_\_\_ \\
deer & 6.000 & 10\% & \_\_\_ \\
dog & 6.000 & 10\% & \_\_\_ \\
frog & 6.000 & 10\% & \_\_\_ \\
horse & 6.000 & 10\% & \_\_\_ \\
ship & 6.000 & 10\% & \_\_\_ \\
truck & 6.000 & 10\% & \_\_\_ \\
\textbf{Total} & \textbf{60.000} & \textbf{100\%} & \textbf{6.000} \\
\end{longtable}

Neste caso:

\[
6000\cdot\frac{6000}{60000}=600
\]

Como todas as classes possuem a mesma quantidade de elementos, a amostra proporcional também fica perfeitamente balanceada.

\begin{center}\rule{0.5\linewidth}{0.5pt}\end{center}

\subsubsection{5. Domicílios de um município}\label{domicuxedlios-de-um-municuxedpio}

\textbf{População:} 20.000 domicílios.
\textbf{Amostra:} 1.000 domicílios.

Estratificação por região do município:

\begin{longtable}[]{@{}lrrr@{}}
\toprule\noalign{}
Estrato & Domicílios & Proporção & Amostra \\
\midrule\noalign{}
\endhead
\bottomrule\noalign{}
\endlastfoot
Centro & 4.000 & 20\% & 200 \\
Norte & 5.000 & 25\% & \_\_\_ \\
Sul & 3.000 & 15\% & \_\_\_ \\
Leste & 6.000 & 30\% & \_\_\_ \\
Oeste & 2.000 & 10\% & \_\_\_ \\
\textbf{Total} & \textbf{20.000} & \textbf{100\%} & \textbf{1.000} \\
\end{longtable}

Por exemplo, para a região Leste:

\[
n_{\text{Leste}} = 1000\cdot\frac{6000}{20000} = 300
\]

O ponto central para mostrar no slide é:

\[
\boxed{\text{A proporção de cada estrato na amostra é igual à sua proporção na população}}
\]

\subsubsection{6. Pacientes de uma clínica --- população balanceada}\label{pacientes-de-uma-cluxednica-populauxe7uxe3o-balanceada}

\textbf{População:} 4.000 pacientes.
\textbf{Amostra:} 400 pacientes.

Estratificação por tipo de atendimento:

\begin{longtable}[]{@{}lrrr@{}}
\toprule\noalign{}
Estrato & População & Proporção & Amostra \\
\midrule\noalign{}
\endhead
\bottomrule\noalign{}
\endlastfoot
Cardiologia & 1.000 & 25\% & 100 \\
Ortopedia & 1.000 & 25\% & \_\_\_ \\
Dermatologia & 1.000 & 25\% & \_\_\_ \\
Oftalmologia & 1.000 & 25\% & \_\_\_ \\
\textbf{Total} & \textbf{4.000} & \textbf{100\%} & \textbf{400} \\
\end{longtable}

Por exemplo:

\[
n_{\text{Cardiologia}} = 400\cdot\frac{1000}{4000} = 100
\]

Como todos os estratos têm o mesmo tamanho, a amostra também fica igualmente distribuída.

\begin{center}\rule{0.5\linewidth}{0.5pt}\end{center}

\subsubsection{7. Produtos de um estoque --- população balanceada}\label{produtos-de-um-estoque-populauxe7uxe3o-balanceada}

\textbf{População:} 10.000 produtos.
\textbf{Amostra:} 1.000 produtos.

Estratificação por categoria:

\begin{longtable}[]{@{}lrrr@{}}
\toprule\noalign{}
Estrato & População & Proporção & Amostra \\
\midrule\noalign{}
\endhead
\bottomrule\noalign{}
\endlastfoot
Eletrônicos & 2.000 & 20\% & 200 \\
Alimentos & 2.000 & 20\% & \_\_\_ \\
Vestuário & 2.000 & 20\% & \_\_\_ \\
Higiene & 2.000 & 20\% & \_\_\_ \\
Papelaria & 2.000 & 20\% & \_\_\_ \\
\textbf{Total} & \textbf{10.000} & \textbf{100\%} & \textbf{1.000} \\
\end{longtable}

Para cada categoria:

\[
1000\cdot\frac{2000}{10000}=200
\]

Esse é outro caso de população perfeitamente balanceada.

\begin{center}\rule{0.5\linewidth}{0.5pt}\end{center}

\subsection{8. Imagens de um dataset de classificação --- população balanceada}\label{imagens-de-um-dataset-de-classificauxe7uxe3o-populauxe7uxe3o-balanceada}

\textbf{População:} 30.000 imagens.
\textbf{Amostra:} 3.000 imagens.

Considere um dataset com 6 classes:

\begin{longtable}[]{@{}lrrr@{}}
\toprule\noalign{}
Estrato --- classe & População & Proporção & Amostra \\
\midrule\noalign{}
\endhead
\bottomrule\noalign{}
\endlastfoot
Cachorro & 5.000 & 16,67\% & 500 \\
Gato & 5.000 & 16,67\% & \_\_\_ \\
Cavalo & 5.000 & 16,67\% & \_\_\_ \\
Pássaro & 5.000 & 16,67\% & \_\_\_ \\
Peixe & 5.000 & 16,67\% & \_\_\_ \\
Coelho & 5.000 & 16,67\% & \_\_\_ \\
\textbf{Total} & \textbf{30.000} & \textbf{100\%} & \textbf{3.000} \\
\end{longtable}

Para a classe ``Cachorro'':

\[
n_h= 3000\cdot\frac{5000}{30000} = 500
\]

Todas as classes possuem a mesma representatividade.

\begin{center}\rule{0.5\linewidth}{0.5pt}\end{center}

\subsection{9. Clientes de um banco --- população não balanceada}\label{clientes-de-um-banco-populauxe7uxe3o-nuxe3o-balanceada}

\textbf{População:} 100.000 clientes.
\textbf{Amostra:} 2.000 clientes.

Estratificação por faixa de renda:

\begin{longtable}[]{@{}lrrr@{}}
\toprule\noalign{}
Estrato & População & Proporção & Amostra \\
\midrule\noalign{}
\endhead
\bottomrule\noalign{}
\endlastfoot
Até R\$ 2 mil & 45.000 & 45\% & 900 \\
R\$ 2 mil a R\$ 5 mil & 30.000 & 30\% & \_\_\_ \\
R\$ 5 mil a R\$ 10 mil & 15.000 & 15\% & \_\_\_ \\
R\$ 10 mil a R\$ 20 mil & 7.000 & 7\% & \_\_\_ \\
Acima de R\$ 20 mil & 3.000 & 3\% & \_\_\_ \\
\textbf{Total} & \textbf{100.000} & \textbf{100\%} & \textbf{2.000} \\
\end{longtable}

Para o primeiro estrato:

\[
n_h = 2000\cdot\frac{45000}{100000} = 900
\]

Aqui os estratos possuem tamanhos bastante diferentes, mas a amostra preserva essas diferenças proporcionalmente.

\begin{center}\rule{0.5\linewidth}{0.5pt}\end{center}

\subsection{10. Dataset de detecção de fraude --- população não balanceada}\label{dataset-de-detecuxe7uxe3o-de-fraude-populauxe7uxe3o-nuxe3o-balanceada}

\textbf{População:} 50.000 transações.
\textbf{Amostra:} 5.000 transações.

Estratificação pelo tipo de transação:

\begin{longtable}[]{@{}lrrr@{}}
\toprule\noalign{}
Estrato & População & Proporção & Amostra \\
\midrule\noalign{}
\endhead
\bottomrule\noalign{}
\endlastfoot
Transações normais & 42.500 & 85\% & 4.250 \\
Fraudes leves & 4.000 & 8\% & \_\_\_\_\_ \\
Fraudes moderadas & 2.000 & 4\% & \_\_\_\_\_ \\
Fraudes graves & 1.000 & 2\% & \_\_\_\_\_ \\
Fraudes críticas & 500 & 1\% & \_\_\_\_\_ \\
\textbf{Total} & \textbf{50.000} & \textbf{100\%} & \textbf{5.000} \\
\end{longtable}

Para as fraudes críticas:

\[
n_{\text{críticas}} = 5000\cdot\frac{500}{50000} = 50
\]

OBS: \textbf{classes raras continuam raras na amostra proporcional}.

Assim, podemos resumir os exemplos 6--10 como:

\begin{longtable}[]{@{}lll@{}}
\toprule\noalign{}
Exemplo & População & Situação \\
\midrule\noalign{}
\endhead
\bottomrule\noalign{}
\endlastfoot
6 & Pacientes por especialidade & Balanceada \\
7 & Produtos por categoria & Balanceada \\
8 & Dataset de imagens & Balanceada \\
9 & Clientes por faixa de renda & Não balanceada \\
10 & Transações financeiras & Não balanceada \\
\end{longtable}

\subsection{TESTES DE MULTIPLA ESCOLHA}\label{testes-de-multipla-escolha-1}

\subsection{Estratos , Proporções e Amostras}\label{estratos-proporuxe7uxf5es-e-amostras}

\begin{longtable}[]{@{}
  >{\raggedright\arraybackslash}p{(\linewidth - 0\tabcolsep) * \real{1.0071}}@{}}
\toprule\noalign{}
\begin{minipage}[b]{\linewidth}\raggedright
\textbf{Questão 5}
\end{minipage} \\
\midrule\noalign{}
\endhead
\bottomrule\noalign{}
\endlastfoot
Uma população é composta por \textbf{40\% de indivíduos da classe A} e \textbf{60\% da classe B}. Deseja-se retirar uma amostra de \textbf{100 indivíduos}.

Qual alternativa descreve corretamente uma \textbf{amostra proporcional} e uma \textbf{amostra balanceada}, respectivamente? \\
\textbf{A)} Proporcional: 50 indivíduos de A e 50 de B; Balanceada: 40 de A e 60 de B. \\
\textbf{B)} Proporcional: 40 indivíduos de A e 60 de B; Balanceada: 50 de A e 50 de B. \\
\textbf{C)} Proporcional: 60 indivíduos de A e 40 de B; Balanceada: 40 de A e 60 de B. \\
\textbf{D)} Proporcional: 40 indivíduos de A e 40 de B; Balanceada: 60 de A e 40 de B. \\
\textbf{E)} Proporcional e balanceada são conceitos equivalentes, portanto ambas devem possuir 50 indivíduos de cada classe. \\
\end{longtable}

\begin{center}\rule{0.5\linewidth}{0.5pt}\end{center}

\subsection{Questão 6 KNN}\label{questuxe3o-6-knn}

\begin{longtable}[]{@{}
  >{\raggedright\arraybackslash}p{(\linewidth - 0\tabcolsep) * \real{1.0065}}@{}}
\toprule\noalign{}
\begin{minipage}[b]{\linewidth}\raggedright
\textbf{Questão 6}
\end{minipage} \\
\midrule\noalign{}
\endhead
\bottomrule\noalign{}
\endlastfoot
Sobre o algoritmo \textbf{K-Nearest Neighbors (KNN)}, assinale a alternativa correta. \\
\textbf{A)} O KNN é um algoritmo paramétrico que precisa determinar obrigatoriamente os parâmetros \(w\) e \(b\). \\
\textbf{B)} O KNN utiliza exclusivamente uma função sigmoide para calcular a classe de uma nova observação. \\
\textbf{C)} O KNN é um algoritmo \textbf{não paramétrico} que pode classificar uma nova observação considerando a distância entre ela e as observações conhecidas. \\
\textbf{D)} No KNN, o valor de \(K\) representa o número de atributos existentes no conjunto de dados. \\
\textbf{E)} O KNN somente pode ser utilizado quando existem exatamente duas classes. \\
\end{longtable}

\begin{center}\rule{0.5\linewidth}{0.5pt}\end{center}

\subsection{Questão 7 - Regressão Logística}\label{questuxe3o-7---regressuxe3o-loguxedstica}

\begin{longtable}[]{@{}
  >{\raggedright\arraybackslash}p{(\linewidth - 0\tabcolsep) * \real{1.0076}}@{}}
\toprule\noalign{}
\begin{minipage}[b]{\linewidth}\raggedright
\textbf{Questão 7 --- Exercício resolvido: Regressão Logística}
\end{minipage} \\
\midrule\noalign{}
\endhead
\bottomrule\noalign{}
\endlastfoot
Considere o modelo de \textbf{Regressão Logística} apresentado em aula: \(\displaystyle z = wX + b\) com:

\(\displaystyle w = 1{,}19\)

e

\(\displaystyle b = -6{,}52\)

Para um aluno que estudou \textbf{6 horas}, temos:

\(\displaystyle z = 1{,}19(6) - 6{,}52\)

\(\displaystyle z = 7{,}14 - 6{,}52\)

\(\displaystyle z = 0{,}62\)

Aplicando a \textbf{função logística (sigmoide)}:

\(\displaystyle P(Y=1 \mid X) = \frac{1}{1+e^{-z}}\)

temos aproximadamente:

\(\displaystyle P(Y=1 \mid X=6) \approx 0{,}65\)

ou seja, aproximadamente \textbf{65\%}.

Considerando um limiar de classificação de \textbf{50\%}, analise as afirmações:

\emph{\textbf{I.} Como a probabilidade calculada é aproximadamente 65\%, o aluno será classificado como \textbf{Aprovado}.}

\emph{\textbf{II.} O valor} \(z=0{,}62\) \emph{corresponde diretamente a uma probabilidade de 62\%.}

\emph{\textbf{III.} A função sigmoide transforma o valor de} \(z\) \emph{em um número entre 0 e 1, que pode ser interpretado como probabilidade.}

Assinale a alternativa \textbf{correta}. \\
\textbf{A)} Apenas I está correta. \\
\textbf{B)} Apenas II está correta. \\
\textbf{C)} Apenas I e II estão corretas. \\
\textbf{D)} Apenas I e III estão corretas. \\
\textbf{E)} I, II e III estão corretas. \\
\end{longtable}

\begin{center}\rule{0.5\linewidth}{0.5pt}\end{center}

\subsection{Questão 8 - KNN}\label{questuxe3o-8---knn}

\begin{longtable}[]{@{}
  >{\raggedright\arraybackslash}p{(\linewidth - 0\tabcolsep) * \real{1.0078}}@{}}
\toprule\noalign{}
\begin{minipage}[b]{\linewidth}\raggedright
\textbf{Questão 8}
\end{minipage} \\
\midrule\noalign{}
\endhead
\bottomrule\noalign{}
\endlastfoot
Considere uma nova observação:

\(\displaystyle X=(4,4)\)

e quatro observações já classificadas:

\textbf{A:} \((2,3)\) --- Azul ; \textbf{B:} \((3,4)\) --- Azul ; \textbf{C:} \((5,4)\) --- Vermelho ; \textbf{D:} \((8,8)\) --- Vermelho ;

Utilizando a \textbf{distância Euclidiana}:

\(\displaystyle d(P,X)=\sqrt{(x_P-x_X)^2+(y_P-y_X)^2}\)

temos:

\(\displaystyle d(A,X)=\sqrt{(2-4)^2+(3-4)^2}=\sqrt{5}\approx 2{,}24\)

\(\displaystyle d(B,X)=\sqrt{(3-4)^2+(4-4)^2}=1\)

\(\displaystyle d(C,X)=\sqrt{(5-4)^2+(4-4)^2}=1\)

\(\displaystyle d(D,X)=\sqrt{(8-4)^2+(8-4)^2}=\sqrt{32}\approx 5{,}66\)

Portanto, considerando \(K=3\), os três vizinhos mais próximos são \textbf{B, C e A}, cujas classes são:

\textbf{Azul, Vermelho, Azul}.

Analise as afirmações:

\textbf{I.} Para \(K=3\), a nova observação será classificada como \textbf{Azul}.

\textbf{II.} O ponto D possui menor distância até \(X\) do que o ponto A.

\textbf{III.} No KNN com \(K=3\), são considerados os três pontos que apresentam as menores distâncias em relação à nova observação.

Assinale a alternativa \textbf{correta}. \\
\textbf{A)} Apenas I está correta. \\
\textbf{B)} Apenas II está correta. \\
\textbf{C)} Apenas I e II estão corretas. \\
\textbf{D)} Apenas II e III estão corretas. \\
\textbf{E)} Apenas I e III estão corretas. \\
\end{longtable}

\begin{center}\rule{0.5\linewidth}{0.5pt}\end{center}

\chapter{Gabarito}\label{gabarito-1}

\begin{longtable}[]{@{}cc@{}}
\toprule\noalign{}
\textbf{Questão} & \textbf{Resposta correta} \\
\midrule\noalign{}
\endhead
\bottomrule\noalign{}
\endlastfoot
\textbf{5} & \textbf{B} \\
\textbf{6} & \textbf{C} \\
\textbf{7} & \textbf{D} \\
\textbf{8} & \textbf{E} \\
\end{longtable}

\chapter{- Análise de Dados: Dilema Viés-Variância e Métricas de Avaliação}\label{anuxe1lise-de-dados-dilema-viuxe9s-variuxe2ncia-e-muxe9tricas-de-avaliauxe7uxe3o}

\section{- Exercícios - Análise de Dados: Dilema Viés-Variância e Métricas de Avaliação}\label{exercuxedcios---anuxe1lise-de-dados-dilema-viuxe9s-variuxe2ncia-e-muxe9tricas-de-avaliauxe7uxe3o}

\chapter{- Relatórios K.P.I. - Avaliação de Modelos e Árvores de Decisão}\label{relatuxf3rios-k.p.i.---avaliauxe7uxe3o-de-modelos-e-uxe1rvores-de-decisuxe3o}

\section{- Exercícios - Relatórios K.P.I. - Avaliação de Modelos e Árvores de Decisão}\label{exercuxedcios---relatuxf3rios-k.p.i.---avaliauxe7uxe3o-de-modelos-e-uxe1rvores-de-decisuxe3o}

\chapter{- Business Intelligence: Random Forest e Ensemble Learning}\label{business-intelligence-random-forest-e-ensemble-learning}

\section{- Exercícios - Business Intelligence: Random Forest e Ensemble Learning}\label{exercuxedcios---business-intelligence-random-forest-e-ensemble-learning}

\chapter{- Big Data e Support Vector Machines}\label{big-data-e-support-vector-machines}

\section{- Exercícios - Big Data e Support Vector Machines}\label{exercuxedcios---big-data-e-support-vector-machines}

\chapter{- Data Driven Decision Making e Redução de Dimensionalidade}\label{data-driven-decision-making-e-reduuxe7uxe3o-de-dimensionalidade}

\section{- Exercícios - Data Driven Decision Making e Redução de Dimensionalidade}\label{exercuxedcios---data-driven-decision-making-e-reduuxe7uxe3o-de-dimensionalidade}

\chapter{- Modelos de Gestão de Risco e Técnicas de Agrupamento}\label{modelos-de-gestuxe3o-de-risco-e-tuxe9cnicas-de-agrupamento}

\section{- Exercícios - Modelos de Gestão de Risco e Técnicas de Agrupamento}\label{exercuxedcios---modelos-de-gestuxe3o-de-risco-e-tuxe9cnicas-de-agrupamento}

\chapter{- Inteligência Artificial, Análise Preditiva e Seleção de Características}\label{inteliguxeancia-artificial-anuxe1lise-preditiva-e-seleuxe7uxe3o-de-caracteruxedsticas}

\section{- Exercícios - Inteligência Artificial, Análise Preditiva e Seleção de Características}\label{exercuxedcios---inteliguxeancia-artificial-anuxe1lise-preditiva-e-seleuxe7uxe3o-de-caracteruxedsticas}

\chapter{- Monitoramento Contínuo e Ciclo de Vida de Modelos de Machine Learning}\label{monitoramento-contuxednuo-e-ciclo-de-vida-de-modelos-de-machine-learning}

\section{- Exercícios - Monitoramento Contínuo e Ciclo de Vida de Modelos de Machine Learning}\label{exercuxedcios---monitoramento-contuxednuo-e-ciclo-de-vida-de-modelos-de-machine-learning}

\bibliography{book.bib}

\end{document}
